\section{Background and Scope}

Rust-for-Linux exposes selected C declarations through bindgen-generated
bindings and hand-written helpers. Safe abstractions then encapsulate unsafe
operations and document invariants such as allocation ownership, error mapping,
lifetime constraints, and allowed execution context. These invariants are often
not visible in a Rust type signature.

The artifact scope is Linux mainline, Rust-enabled builds, and the
Rust-for-Linux surfaces that connect generated bindings, helpers, and
\texttt{rust/kernel} abstractions. The main labeled evaluation is x86\_64. An
arm64 replay slice is used for external-validity evidence, but it does not add
independent labels. The paper therefore avoids claims about all architectures,
all Rust kernel branches, or complete API evolution coverage.

Labels follow the repository review guide. \truewrapper{} is not counted as
\truesemantic{}. \truesemantic{} means that C-side drift, Rust exposure, and a
plausible stale-contract dependency are supported by evidence; it does not mean
a concrete runtime bug has been demonstrated.
